\subsection{BatchReceive机制与批内参与方互异条件}
\label{subsec:batchreceive-condition}

\kwaay 由KEM、数字签名和split-KEM等密码组件构成，其完整握手和计算安全证明不在本文分析范围内\cite{collins2024kwaay}。与批次问题直接相关的是接收侧接口：发送方调用\sendaction 产生协议消息；原文第4.1节说明，当多条发送方消息对应同一接收方临时状态时，接收方以该共享状态调用\batchreceive\cite{collins2024kwaayfull}。一次调用的输入可概括为
\[
  S=\bigl((\mathit{pk}_{j_1},\mathit{prek}_{j_1},m_{j_1}),\ldots,
           (\mathit{pk}_{j_d},\mathit{prek}_{j_d},m_{j_d})\bigr),
  \qquad d\geq 1,
\]
其中，每个分量包含所声称发送方的公钥、预密钥束和协议消息。该向量只表示一次接收操作会组合多个发送方关联输入，不涉及各分量的解封装和密钥派生过程。

原文第5.1节（第21页）要求同次\batchreceive 调用中的元素对应不同参与方\cite{collins2024kwaayfull}。该要求约束同一批次内元素之间的关系，并不要求某个参与方在整个执行中只出现一次，也不限制其只能产生一条消息。规范明确了目标关系，但未在该处指定由调用者、批次构造组件还是准入层负责维护。

分析对象是候选条目完成组合之后、进入后续处理之前的边界，下文称为批次准入。“批次准入”仅概括这一分析边界，并非为K-Waay增加新算法。该边界需要刻画同次调用中各条目的参与方投影及其相互关系，而不是重新描述完整密码过程。

\subsection{比较维度与问题实例}
\label{subsec:identity-coordinates}

为区分不同限制实际比较的对象，将参与方、发送实例和消息三个分量分别表示。批次准入所需的条目信息抽象为
\begin{equation}
  E_i=(A_i,\oid_i,m_i),
  \label{eq:abstract-entry}
\end{equation}
其中，$A_i$是参与方分量，$\oid_i$是发送实例标识，$m_i$是消息分量。中文稿中的$\oid$对应Tamarin源模型的$\mathit{sid}$，表示一次\texttt{SendMessage}规则实例化产生的发送实例，不是完整协议的会话标识、transcript或事件时间点。式\eqref{eq:abstract-entry}仅用于分析，不表示协议线路编码。

\begin{table}[htbp]
  \centering
  \caption{身份坐标与模型映射}
  \label{tab:identity-coordinate-map}
  \footnotesize
  \begin{tabularx}{\linewidth}{@{}L{0.13\linewidth}C{0.13\linewidth}C{0.15\linewidth}YY@{}}
    \toprule
    概念 & 论文记号 & Tamarin记号 & 作用 & 解释边界 \\
    \midrule
    参与方 & $A$ & $A$ & 标识建模协议主体 & 不自动等同于账户或密钥编码 \\
    发送发生 & $\oid$ & $\mathit{sid}$ & 区分发送方产生条目的发生 & 发生不同不推出参与方不同 \\
    消息 & $m$ & $m$ & 表示发送方产生的符号消息 & 消息不同不推出参与方不同 \\
    槽位 & $i$ & 线性收集状态 & 标识批次内位置 & 位置不同不保证条目内容互异 \\
    批次/接收上下文 & $(\mathit{bid},\mathit{rst})$ & $(\mathit{bid},\mathit{rst})$ & 限定共同处理作用域 & 不表示完整接收方状态 \\
    \bottomrule
  \end{tabularx}
\end{table}



表\ref{tab:identity-coordinate-map}中的维度具有不同作用：槽位和$(\mathit{bid},\mathit{rst})$限定条目所在的批次上下文，消息分量区分符号消息，发送实例标识区分具体发送，参与方分量则是目标条件比较的对象。$A$表示符号模型中的协议主体，不自动等同于账户、公钥字节串或数据库记录；具体系统仍需说明实现对象到参与方投影的映射。

考虑候选条目$(A,\oid_1,m_1)$和$(A,\oid_2,m_2)$，其中$\oid_1\neq\oid_2$且$m_1\neq m_2$。二者来自不同发送实例并携带不同消息，但参与方投影相同。这一最小实例说明，消息或发送实例互异在静态关系上不能保证参与方互异；其是否构成模型中的实际行为，还需通过共同生命周期的可达性验证。

\subsection{研究问题}
\label{subsec:research-questions}

基于上述区分，提出三个研究问题：
\begin{enumerate}[label=\textbf{RQ\arabic*：},leftmargin=4.2em,labelsep=0.5em]
  \item 移除批内参与方互异条件后，重复参与方条目能否进入同一批次并产生接受事件？
  \item 一个匹配发送来源能否在同一批次和接收方上下文中对应两个\receiveraccept 事件？
  \item 消息互异约束能否替代参与方互异条件，还是同一参与方的不同消息仍可同时获准进入批次？
\end{enumerate}

为使验证差异能够归因于准入条件，三种模型保持发送、收集、来源匹配和顺序处理过程不变，仅调整准入阶段比较的维度。

